% ---------------------------------------------------------------------------
% Glossaire et liste des sigles
% ---------------------------------------------------------------------------

\frontchapter{Glossaire et liste des sigles}

\renewcommand{\arraystretch}{1.25}

% =============================================================================
% LISTE DES SIGLES
% =============================================================================

\section*{Liste des sigles}
\phantomsection
\addcontentsline{toc}{section}{Liste des sigles}

\small

\begin{longtable}{
    @{}
    >{\bfseries\raggedright\arraybackslash}p{0.18\textwidth}
    >{\raggedright\arraybackslash}p{0.76\textwidth}
    @{}
}

\toprule
\TableHeader
Sigle & Signification \\
\midrule
\endfirsthead

\toprule
\TableHeader
Sigle & Signification \\
\midrule
\endhead

\bottomrule
\endlastfoot

API
&
\textit{Application Programming Interface}. Interface permettant à plusieurs
composants logiciels ou applications de communiquer entre eux.
\\

CAO
&
Conception assistée par ordinateur. Domaine regroupant les outils et formats de
dessin technique, comme les fichiers \texttt{DWG} ou \texttt{DXF}.
\\

CRS
&
\textit{Coordinate Reference System}. Système de référence de coordonnées
utilisé pour localiser des objets géographiques.
\\

DGN
&
Format de fichier CAO notamment associé à certains outils de dessin technique.
\\

DWG
&
Format de fichier CAO couramment utilisé par les logiciels de dessin assisté par
ordinateur.
\\

DXF
&
\textit{Drawing Exchange Format}. Format CAO utilisé pour échanger des dessins
entre logiciels.
\\

EPSG
&
Référentiel de codes permettant d'identifier les systèmes de coordonnées et
les transformations géodésiques.
\\

FME
&
\textit{Feature Manipulation Engine}. Plateforme d'intégration et de
transformation de données développée par Safe Software.
\\

GDAL
&
\textit{Geospatial Data Abstraction Library}. Bibliothèque libre permettant de
lire, écrire et transformer de nombreux formats de données géospatiales.
\\

LAS
&
Format de fichier utilisé pour stocker des nuages de points, notamment issus de
relevés LiDAR.
\\

LAZ
&
Version compressée du format \texttt{LAS}, utilisée pour réduire le volume de
stockage des nuages de points.
\\

OGR
&
Composant de GDAL principalement consacré à la lecture, à l'écriture et au
traitement des données vectorielles.
\\

PDAL
&
\textit{Point Data Abstraction Library}. Bibliothèque spécialisée dans la
lecture et le traitement des nuages de points.
\\

PR
&
\textit{Pull Request}. Demande d'intégration de modifications dans un projet
versionné, généralement relue avant d'être acceptée.
\\

RSE
&
Responsabilité sociétale des entreprises.
\\

SDF
&
Format de données géographiques ou CAO selon les contextes, étudié puis écarté
du périmètre initial de développement.
\\

SGBD
&
Système de gestion de base de données.
\\

SIG
&
Système d'information géographique.
\\

SQL
&
\textit{Structured Query Language}. Langage utilisé pour interroger et
manipuler des bases de données relationnelles.
\\

\end{longtable}

\normalsize

\vspace{0.6cm}

% =============================================================================
% GLOSSAIRE
% =============================================================================

\section*{Glossaire}
\phantomsection
\addcontentsline{toc}{section}{Glossaire}

\small

\begin{longtable}{
    @{}
    >{\bfseries\raggedright\arraybackslash}p{0.24\textwidth}
    >{\raggedright\arraybackslash}p{0.70\textwidth}
    @{}
}

\toprule
\TableHeader
Terme & Définition \\
\midrule
\endfirsthead

\toprule
\TableHeader
Terme & Définition \\
\midrule
\endhead

\bottomrule
\endlastfoot

ArcPy
&
Bibliothèque Python fournie par Esri permettant d'automatiser des traitements
réalisés avec ArcGIS et d'accéder aux propriétés des données géographiques
Esri.
\\

Base de données spatiale
&
Base de données capable de stocker, d'interroger et de manipuler des objets
géographiques tels que des points, des lignes et des polygones.
\\

Conda
&
Gestionnaire d'environnements et de dépendances Python. Dans le projet, il a été
retenu pour construire un environnement de packaging plus stable autour de GDAL
et PyInstaller.
\\

Donnée raster
&
Donnée géographique organisée sous la forme d'une grille de cellules ou de
pixels, utilisée notamment pour représenter des images aériennes, des modèles
numériques ou des cartes continues.
\\

Donnée vectorielle
&
Donnée géographique représentée par des géométries telles que des points, des
lignes ou des polygones, auxquelles peuvent être associés des attributs.
\\

Documentation des métadonnées
&
Étape consistant à extraire les informations décrivant une ressource : nom,
format, projection, géométrie, attributs, emprise ou encore dates disponibles.
Dans Scan, cette étape est aussi appelée \textit{lookup}.
\\

Docker
&
Outil permettant d'exécuter une application dans un environnement isolé et
reproductible. Il a été utilisé pour générer l'exécutable Windows dans des
conditions maîtrisées.
\\

Driver
&
Composant logiciel permettant à GDAL/OGR de lire ou d'écrire un format de
fichier ou une source de données particulière.
\\

Emprise spatiale
&
Rectangle minimal décrivant les limites géographiques d'un jeu de données. Elle
est généralement définie par ses coordonnées minimales et maximales.
\\

Enveloppe convexe
&
Plus petit polygone convexe contenant l'ensemble des géométries d'un jeu de
données.
\\

GeoPackage
&
Format de fichier géographique pouvant contenir plusieurs couches dans un même
fichier, par exemple des données vectorielles, des rasters ou des tables.
\\

GeoJSON
&
Format d'échange de données géographiques textuel, couramment utilisé pour
représenter des objets vectoriels et leurs attributs.
\\

GeoParquet
&
Format de stockage géographique reposant sur Parquet. Il permet de stocker des
données vectorielles de manière efficace, notamment pour des usages d'analyse.
\\

Géodatabase
&
Structure de stockage de données géographiques utilisée notamment par les
solutions Esri. Elle peut être locale ou hébergée dans un système de gestion de
base de données.
\\

Géomatique
&
Discipline regroupant les méthodes et technologies utilisées pour acquérir,
stocker, traiter, analyser et représenter des données localisées.
\\

Métadonnée
&
Information décrivant une ressource de données, par exemple son nom, sa
structure, son format, sa projection, son emprise ou sa date de modification.
\\

Laspy
&
Bibliothèque Python utilisée pour lire et écrire des fichiers de nuages de
points aux formats \texttt{LAS} et \texttt{LAZ}.
\\

LibreDWG
&
Bibliothèque libre spécialisée dans la lecture des fichiers \texttt{DWG}. Elle a
été étudiée pour accéder au contenu des données CAO dans le cadre de Scan.
\\

Listing
&
Étape de Scan consistant à inventorier les ressources accessibles dans une
source de données avant de les documenter ou de les signer.
\\

Lookup
&
Nom utilisé dans la documentation technique de Scan pour désigner l'étape de
documentation. Elle consiste à extraire les métadonnées décrivant une ressource.
\\

MkDocs
&
Outil permettant de générer un site statique de documentation. Il a été utilisé
avec \textit{mkdocstrings} pour produire une documentation à partir du code
Python.
\\

Nuage de points
&
Ensemble de points décrivant une surface ou un objet en trois dimensions. Ces
données peuvent être très volumineuses et sont souvent stockées dans des
fichiers \texttt{LAS} ou \texttt{LAZ}.
\\

Oracle Spatial
&
Extension spatiale d'Oracle permettant de stocker et d'interroger des données
géographiques dans une base Oracle.
\\

Packaging
&
Processus consistant à regrouper une application et ses dépendances afin de
permettre son installation ou son exécution dans un environnement cible.
\\

PostGIS
&
Extension spatiale du système de gestion de base de données PostgreSQL. Elle
permet de stocker et d'interroger des géométries ainsi que d'effectuer des
traitements spatiaux.
\\

PyInstaller
&
Outil permettant de transformer une application Python et ses dépendances en
un exécutable pouvant être lancé sans installation préalable de Python.
\\

Scan
&
Produit d'ISOGEO chargé d'explorer des sources de données, d'identifier les
ressources qu'elles contiennent et d'extraire les métadonnées nécessaires à leur
documentation dans l'écosystème ISOGEO.
\\

Signature
&
Empreinte calculée à partir d'une ressource afin de vérifier son intégrité et de
détecter une éventuelle évolution de son contenu.
\\

Shapefile
&
Format de fichier vectoriel Esri très utilisé pour échanger des données
géographiques.
\\

Source de données
&
Emplacement ou système à partir duquel Scan analyse des ressources. Une source
peut notamment correspondre à une base de données, une géodatabase, un dossier
ou un fichier géographique.
\\

SQL Server
&
Système de gestion de base de données relationnelle de Microsoft, pouvant aussi
stocker des données spatiales.
\\

Test de régression
&
Test visant à vérifier qu'une modification du logiciel n'altère pas un
comportement précédemment validé. Dans ce projet, les résultats Python sont
comparés aux résultats produits par les traitements FME de référence.
\\

Workbench FME
&
Fichier de travail FME décrivant graphiquement une chaîne de lecture, de
transformation et d'écriture de données.
\\

\end{longtable}

\normalsize
